\section{Discussion}
\label{sec:discussion}

The clearest result is the importance of learned, transferred
representations.  Fixed colour/texture/gradient summaries provide a valid
baseline but lose spatial and semantic structure; SIFT-BoVW adds local
invariance without closing this gap.  Scratch ResNet18 learns far stronger
features from 40 examples per class, yet pretrained ResNet18 more than doubles
its top-1 accuracy.  EfficientNet-B0 then improves accuracy with fewer
parameters than ResNet18, suggesting that its architecture and pretrained
representation are well matched to the available compute and data.  Because
the scratch/pretrained schedules differ, this evidence is consistent with a
pretraining benefit but does not isolate that benefit causally.

The traditional rewrite also sharpens what the classical comparison means.
Joint colour statistics, multi-scale texture, HOG, and spatial soft-assignment
BoVW are complementary enough that the 3,704-D LinearSVC is the primary
Stage~2 traditional model.  The Stage~4 Random Forest is a different,
2,424-D, 300-tree handcrafted-only comparator without BoVW.  It remains useful
for a fixed corruption study because that evaluator can deterministically
recompute its per-image input and reuse the fitted normaliser, whereas it does
not reconstruct the training-fitted visual vocabulary.  Its robustness curve
must therefore not be read as the robustness of the best-scoring traditional
model.

The strict augmentation comparison is more modest but actionable.  Strong
augmentation improves the pretrained ResNet18 validation score without
altering its initial weights, sampling order, optimiser, or stopping policy.
The result is consistent with nuisance variability in iNaturalist images and
supports the final use of colour, affine, crop, and erasing perturbations.
However, the clean improvement does not make the model invariant to severe
degradation.  Motion blur removes high-frequency edges and local markings that
may be diagnostic, plausibly contributing to its large top-1 fall.  Training
with blur or evaluating corruption-aware
augmentation would be natural future work, but doing so after observing the
held-out corruption scores would require a new validation protocol and test
set.

Confusion and Grad-CAM analyses reveal two distinct error regimes.  Some
fine-grained pairs are correctly localised yet remain inseparable, indicating
that global $224\!\times\!224$ classification may discard diagnostic detail.
Other high-confidence failures align with conspicuous colour or texture, which
may be predictive in training but not species-specific.  A promising extension
would combine higher-resolution or part-based attention with taxonomic
hierarchy, while explicitly testing whether explanations remain stable under
background and corruption interventions.

Several limitations bound our conclusions.  The balanced 500-class subset is
reproducible but does not measure full 10,000-class or long-tail performance.
Forty training images per class constrain scratch learning, and ten test images
per class make per-class metrics noisy.  We evaluate one seed and a limited
hyperparameter budget on one M4 Pro, so small model differences may not be
statistically stable.  Traditional cached-feature timing and deep forward
timing have different scopes.  The corruption suite covers four synthetic
degradations, not distribution shifts in geography, season, camera, or
observer behaviour.  Finally, Grad-CAM is coarse and gradient-based; its maps
are useful for forming hypotheses, not certifying faithful reasoning.  A
further provenance limitation is retrospective: the integrated repository
contains validation and test evidence for rewritten Stage~2, but no surviving
artifact independently proves that the selected configuration was immutably
locked before the first test access.  The restored selection manifest freezes
the current hashes, allowed models, and validation grid going forward; it does
not manufacture a historical chronology.

Despite these limitations, the project has two practical strengths.  First,
the current split, model, prediction, metric, and selection artifacts are
cross-hashed, so subsequent drift or unauthorised test access is detectable.
Second,
the combination of clean metrics, macro scores, classwise errors, corruptions,
and explanations gives a more complete account than accuracy alone.  The
best model is genuinely useful on clean images, but the robustness and failure
analyses show exactly why it should not yet be treated as a dependable field
identifier without uncertainty handling and expert review.
